% =============================================================
%  CAPÍTULO II — MARCO TEÓRICO DEL PROYECTO TECNOLÓGICO
% =============================================================

\section{Marco conceptual}

En el contexto de las evaluaciones por competencia curricular, los métodos de evaluación juegan un papel crucial. Estos métodos son los enfoques y herramientas utilizados para medir el grado en que los estudiantes han adquirido las competencias específicas delineadas en el plan de estudios. Pueden incluir desde pruebas tradicionales hasta proyectos prácticos, pasando por evaluaciones basadas en la resolución de problemas o en la presentación oral, entre otros.

La competencia, en este contexto, se refiere a la capacidad demostrada por un estudiante para aplicar conocimientos, habilidades y actitudes de manera efectiva y apropiada en situaciones específicas. En lugar de simplemente memorizar información, la competencia implica la capacidad de utilizar ese conocimiento de manera significativa y transferible en el mundo real.

El aprendizaje individualizado es un enfoque educativo que reconoce las diferencias individuales entre los estudiantes y adapta el proceso de enseñanza y evaluación para satisfacer las necesidades únicas de cada uno. En el contexto de las evaluaciones por competencia curricular, esto podría implicar la personalización de las evaluaciones para tener en cuenta los estilos de aprendizaje, los intereses y las habilidades individuales de cada estudiante.

La retroalimentación significativa se refiere a los comentarios proporcionados a los estudiantes después de una evaluación con el propósito de guiar su aprendizaje de manera efectiva. En el contexto de las evaluaciones por competencia curricular, la retroalimentación significativa se centra en destacar no solo los errores, sino también los logros del estudiante, identificando áreas de mejora y proporcionando orientación concreta sobre cómo avanzar.

La discapacidad intelectual es una condición que afecta la capacidad de una persona para comprender, procesar y retener información, así como para participar en actividades académicas de la misma manera que sus pares sin discapacidad. En el contexto de las evaluaciones por competencia curricular, es importante considerar las necesidades y adaptaciones específicas que puedan requerir los estudiantes con discapacidades intelectuales para participar de manera equitativa en el proceso de evaluación.

El currículum, también conocido como plan de estudios, se refiere al conjunto de objetivos educativos, contenidos, métodos de enseñanza y criterios de evaluación diseñados para guiar el aprendizaje en un programa educativo específico. En el contexto de las evaluaciones por competencia curricular, el currículum proporciona el marco sobre el cual se basan las evaluaciones, delineando las competencias que se espera que los estudiantes adquieran y demuestren durante su proceso educativo.

\subsection{Evolución de las evaluaciones educativas}

El concepto de evaluación ha ido evolucionando en consonancia con el concepto de educación predominante. Es así que, desde una evaluación centrada en el acto de juzgar el valor de las cosas, se ha evolucionado hacia una evaluación que pretendía asignar valores precisos de medición a determinados objetos educativos.

En la década de los 30 Ralph Tyler inicia un movimiento de la evaluación en función del logro de determinados objetivos formulados con antelación, produciendo un cambio relevante para poder comprender el porqué del proceso educativo, el cual siempre estaba enfocado a ver los resultados del aprendizaje.

En los años 60 la evaluación se redefine como una etapa de ``recopilación y uso de información para tomar decisiones sobre un programa educativo''. Este enfoque transforma la percepción previa de la evaluación como un mero instrumento de control y medición para la valoración final de un proceso. En lugar de eso, se concibe como un medio que proporciona la oportunidad de retroalimentar continuamente el proceso educativo y respaldar la toma de decisiones.

A inicios del siglo XXI se tiende a adoptar una concepción ecléctica del proceso de evaluación, definiéndola como el proceso de delinear, obtener, procesar y proporcionar información válida, confiable y oportuna que nos permita juzgar el mérito o valía de programas, procedimientos y productos con el fin de tomar decisiones.

En la actualidad se considera que es crucial evaluar la comprensión del estudiante respecto a la actividad a realizar, incluyendo su significado, sentido, plenitud y la manera en que se alcanza esa comprensión.

\subsection{Evaluación por competencias}

La evaluación por competencias implica la recopilación de evidencias mediante actividades de aprendizaje y la formulación de valoraciones sobre la medida y la naturaleza del progreso del estudiante en relación con los resultados de aprendizaje esperados.

La evaluación por competencias ofrece nuevas oportunidades a los estudiantes al generar entornos significativos de aprendizaje que acercan sus experiencias académicas al mundo profesional, y donde pueden desarrollar una serie de capacidades integradas y orientadas a la acción, con el objetivo de ser capaces de resolver problemas prácticos o enfrentarse a situaciones ``auténticas''. Estas competencias están compuestas por un conjunto de estructuras de conocimiento, así como habilidades cognitivas, interactivas y afectivas, actitudes y valores, que son necesarias para la ejecución de tareas, la solución de problemas y un desempeño eficaz en una determinada profesión, organización, posición o rol.

\subsection{Evaluación por competencias en matemáticas}

La competencia en matemáticas es un tema de gran relevancia en los países que han adoptado el enfoque competencial para la enseñanza de la materia. Este enfoque destaca la necesidad de reemplazar un currículum centrado únicamente en la adquisición de contenidos para lograr el éxito académico. En su lugar, se aboga por un currículum que se centre en la alfabetización matemática de los estudiantes, permitiéndoles utilizar de manera integral y efectiva los conocimientos matemáticos en diversas situaciones cotidianas en las que estos son necesarios.

Desde esa posición, entendemos la competencia matemática como aquella que nos permite dar un uso funcional al conocimiento matemático en diversas situaciones desde una profunda comprensión. Los estudiantes, cuando se enfrentan a problemas en contextos del mundo real, tendrán que activar las competencias matemáticas pertinentes para resolver el problema. Para ello necesitan disponer de una cultura matemática que les permita analizar, razonar y comunicar efectivamente, al plantear, resolver e interpretar problemas de naturaleza matemática.

El desarrollo de la competencia matemática no es un proceso de aprendizaje con un punto final. Nunca podemos afirmar que, después de un determinado proceso educativo, se haya adquirido completamente la competencia en cuestión. Es un proceso continuo; cada vez que se enfrentan nuevas situaciones y se relacionan nuevos conocimientos, se aumenta la competencia.

\subsection{Tipos de preguntas en evaluaciones por competencias}

En la evaluación por competencias, el diseño de preguntas juega un papel crucial para medir de manera efectiva las habilidades y conocimientos específicos que se buscan evaluar. Estas preguntas están diseñadas para ir más allá de la simple memorización de información, buscando evidenciar la aplicación práctica de conocimientos y habilidades en situaciones del mundo real. A continuación, se muestran varios tipos de preguntas comunes:

\begin{itemize}[leftmargin=*]
  \item \textbf{Preguntas de Aplicación:} Desafían a los evaluados a aplicar sus conocimientos y habilidades en situaciones prácticas. Por ejemplo, ``Usando tus juguetes, muéstrame cómo harías un castillo alto con bloques. ¿Puedes explicar qué partes usaste y por qué?''
  \item \textbf{Estudios de Caso:} Presentan un escenario detallado y piden a los evaluados que analicen, interpreten y propongan soluciones basadas en sus conocimientos.
  \item \textbf{Problemas de Resolución:} Requieren que los evaluados resuelvan problemas utilizando sus habilidades y conocimientos.
  \item \textbf{Preguntas Situacionales:} Presentan situaciones hipotéticas o reales y piden a los evaluados que tomen decisiones o proporcionen soluciones basadas en su conocimiento y experiencia.
\end{itemize}

Estos tipos de preguntas están diseñados para evaluar no solo el conocimiento teórico, sino también la capacidad de aplicar ese conocimiento en contextos prácticos. Además, permiten medir habilidades como el pensamiento crítico, la resolución de problemas, la toma de decisiones y la capacidad de trabajar en situaciones del mundo real.

\subsection{Integración de diferentes tipos de preguntas}

La integración de diversos tipos de preguntas en una evaluación es fundamental para capturar de manera integral las habilidades y conocimientos de los estudiantes. Al combinar preguntas de aplicación, estudios de caso, problemas de resolución y preguntas situacionales, se abarcan diferentes niveles de comprensión y se evalúan habilidades específicas, desde la aplicación práctica hasta el pensamiento crítico. Esta diversidad no solo se adapta a diferentes estilos de aprendizaje, sino que también proporciona a los educadores una perspectiva más completa del rendimiento estudiantil.

La importancia radica en la capacidad de estas preguntas para reflejar la complejidad del mundo real y preparar a los estudiantes para enfrentar desafíos diversos. Además, motiva la participación activa al mantener la evaluación interesante y relevante. En última instancia, la integración de distintos tipos de preguntas en una evaluación no solo mide el conocimiento superficial, sino que impulsa a los estudiantes a aplicar, analizar y tomar decisiones, promoviendo así un aprendizaje más significativo.

\subsection{Uso de elementos multimedia en preguntas educativas}

La inclusión de elementos multimedia en las evaluaciones dirigidas a niños desempeña un papel crucial en fomentar un ambiente de aprendizaje estimulante y accesible. La incorporación de imágenes coloridas, videos interactivos y elementos visuales no solo captura la atención de los niños, sino que también facilita la comprensión de conceptos y situaciones, promoviendo una experiencia de evaluación más emocionante y participativa.

Los elementos multimedia no solo son atractivos, sino que también se alinean con la naturaleza lúdica del aprendizaje infantil. Utilizar imágenes y videos puede convertir la evaluación en una experiencia más amigable y divertida para los niños, lo que contribuye a un ambiente de aprendizaje positivo. Además, al simular escenarios de la vida real a través de formatos multimedia, se evalúa la capacidad de los niños para aplicar sus habilidades en contextos prácticos de manera más intuitiva.

La claridad y la comprensión se mejoran significativamente mediante elementos multimedia adaptados a la edad. Los niños pueden beneficiarse al ver visualmente conceptos abstractos o seguir instrucciones a través de ilustraciones, lo que les permite expresar sus conocimientos y habilidades de manera más efectiva.

\subsection{La tecnología educativa}

La tecnología educativa se ocupa del análisis de los medios, materiales, sitios web y plataformas tecnológicas utilizados en los procesos de aprendizaje. Dentro de su ámbito, se incluyen recursos diseñados inicialmente para abordar las necesidades de los usuarios en términos formativos e instructivos.

En cuanto al uso de tecnologías en la educación, la conexión entre educación y tecnología no es algo novedoso, sino más bien un elemento constante a lo largo de la historia. En la actualidad, se enfatiza que no se trata simplemente de aumentar la intensidad del uso de la tecnología, sino más bien de comprender claramente los beneficios que las opciones tecnológicas pueden aportar.

Actualmente la tecnología facilita estrategias didácticas enriquecedoras, permitiendo diseños atractivos y motivadores para alcanzar objetivos educativos. Además, propicia un cambio de roles tanto para docentes como para estudiantes al aumentar la persistencia, focalización y accesibilidad en las tareas, posibilitando un mayor acompañamiento docente en el proceso de enseñanza-aprendizaje.

\subsection{Aplicaciones educativas}

Las aplicaciones educativas representan valiosas herramientas diseñadas para consolidar los conocimientos de niños en diversas áreas de aprendizaje. Estas aplicaciones, enriquecidas con imágenes, sonidos, dibujos y animaciones adaptadas a las distintas edades y disciplinas, se erigen como elementos cruciales para el proceso educativo. Se reconoce que la inclusión de elementos visuales y auditivos facilita el procesamiento de información por parte de los seres humanos.

La relevancia de estas aplicaciones en el aprendizaje de los estudiantes radica en su capacidad para ofrecer un entorno atractivo, estimulante y efectivo. Se convierten así en factores significativos que podrían potenciar el desarrollo educativo de los niños. Es imperativo que estas aplicaciones educativas se incorporen progresivamente en diversas instituciones educativas, especialmente considerando su sinergia con la omnipresente herramienta de internet y el acceso a través de tabletas.

Además de ser instrumentos para el aprendizaje individual, estas aplicaciones desempeñan un papel esencial en la promoción de la colaboración y la mejora de la comunicación entre los estudiantes. Facilitan el compartir experiencias y el trabajo en equipo, contribuyendo así a un ambiente educativo más interactivo y enriquecedor.

\subsection{Desarrollo de aplicaciones para evaluación educativa}

La importancia de las aplicaciones educativas radica en su capacidad para proporcionar experiencias de aprendizaje interactivas y personalizadas. Al incorporar elementos visuales, auditivos y táctiles, estas aplicaciones estimulan la participación activa de los estudiantes, fomentando un ambiente educativo más dinámico y atractivo. Además, ofrecen la flexibilidad necesaria para adaptarse a diversos estilos de aprendizaje, permitiendo un enfoque más individualizado y centrado en el estudiante.

La selección adecuada de tecnologías para el desarrollo de aplicaciones educativas es esencial para garantizar su eficacia y alineación con los objetivos pedagógicos. Es crucial que estas tecnologías sean intuitivas, seguras y adaptables a diferentes dispositivos, asegurando así un acceso fácil y equitativo para todos los estudiantes.

Para el desarrollo de esta aplicación centrada en la evaluación por competencias curriculares, se han seleccionado tecnologías avanzadas que contribuirán significativamente a la eficacia y versatilidad del proyecto. Angular, reconocido por su robustez en la creación de interfaces de usuario interactivas, aporta una experiencia visual envolvente para los estudiantes. NestJS centraliza la lógica del backend y los protocolos de seguridad, mientras que PostgreSQL garantiza un almacenamiento estructurado y escalable.

\subsection{Programación Extrema (XP)}

La Programación Extrema (XP) es una metodología ágil de desarrollo de software que se enfoca en la entrega rápida de valor al cliente mediante iteraciones cortas y prácticas técnicas rigurosas. XP es independiente de herramientas o lenguajes particulares, brindando a los equipos la flexibilidad para construir software de alta calidad en entornos de alta incertidumbre.

Las prácticas de esta metodología se alinean con un ciclo de vida iterativo y se fundamentan en valores como la comunicación constante, la simplicidad del diseño y el coraje para refactorizar y mejorar el sistema continuamente. El objetivo de XP es entregar valor al cliente a través de lanzamientos frecuentes de software funcional, reduciendo los riesgos y optimizando el proceso de desarrollo mediante la colaboración y la excelencia técnica.

A continuación, se describen las prácticas fundamentales de la Programación Extrema, organizadas según las cuatro actividades principales del ciclo de desarrollo propuesto por Pressman:

\begin{itemize}[leftmargin=*]
  \item \textbf{Fase de Planeación:} Esta etapa inicial establece la hoja de ruta del proyecto mediante una colaboración intensa entre el equipo de desarrollo y el cliente. A través del Juego de Planificación, el cliente presenta las necesidades en forma de Historias de Usuario, las cuales son discutidas, estimadas en esfuerzo por el equipo técnico y priorizadas según su valor.
  \item \textbf{Fase de Diseño:} La prioridad en XP es mantener el Diseño Simple, buscando siempre la solución más sencilla que cumpla con los requisitos actuales. Se evita la complejidad innecesaria y la anticipación de funcionalidades futuras.
  \item \textbf{Fase de Codificación:} El Desarrollo Dirigido por Pruebas (TDD) es una práctica fundamental, donde se escriben pruebas automatizadas que fallan antes de escribir el código funcional. La Refactorización se realiza de forma constante para mejorar la estructura interna del código sin alterar su comportamiento externo.
  \item \textbf{Fase de Prueba y Entrega:} Los Lanzamientos Pequeños son una práctica clave, donde se liberan versiones funcionales del sistema al cliente con mucha frecuencia. Esto permite obtener retroalimentación temprana sobre el software real.
\end{itemize}

\section{Marco referencial}

La evaluación basada en competencias ha experimentado una integración cada vez mayor con la tecnología en los últimos años, dando lugar al desarrollo de diversas aplicaciones destinadas a facilitar y mejorar este proceso. A continuación, se presentan algunas de estas aplicaciones:

\textbf{SCALA (Enfoque Analítico de Aprendizaje Escalable):} Se trata de una herramienta que genera análisis de aprendizaje a partir de datos recopilados de diversos repositorios educativos en 2015. SCALA elabora rúbricas detalladas considerando el modelo curricular basado en competencias de la Universidad de Deusto y los datos extraídos de los repositorios educativos sobre las actividades realizadas por los estudiantes durante los cursos.

\textbf{CAT (Herramienta de Evaluación de Competencias):} Es una plataforma que permite la evaluación de competencias genéricas a nivel universitario utilizando blogs del año 2016. Su objetivo es mejorar los procesos de autorreflexión y proporcionar retroalimentación junto con rúbricas para evaluar ciertos indicadores de éxito asociados a competencias específicas.

\textbf{SECAT (Herramienta de Evaluación de Competencias en Ingeniería de Software):} Se trata de un marco de evaluación por competencias para la ingeniería de software del año 2015, que adapta y enriquece el enfoque de Rauner. Este enfoque utiliza evaluaciones individuales y de equipo a través de cuestionarios cuyas preguntas están relacionadas con diversas dimensiones de competencia.

\textbf{MCAT (Herramienta de Evaluación de Competencias en Obstetricia):} Adaptada de las competencias esenciales de la Confederación Internacional de Matronas (ICM) y del Colegio Americano de Enfermeras Parteras (ACNM) en 2018, esta herramienta se basa en la teoría del principiante al experto para determinar niveles de competencia.

\textbf{LACAMOLC:} Propuesta de plataforma web de 2014 que ofrece a docentes y estudiantes un panel de control que recopila datos sociales y de uso de diversas tecnologías de aprendizaje, como Moodle y Google Apps. Utiliza Pentaho como herramienta de análisis.

\textbf{AEEA (Arquitectura de Evaluación de Motor Adaptativo):} Suite de herramientas web adaptativas de 2013 que permite diseñar, evaluar y desplegar análisis de aprendizaje para cursos en línea y mixtos de instituciones de educación superior bajo el marco europeo de cualificaciones (EQF).

Estas aplicaciones han surgido como soluciones que intentan abordar varios aspectos clave, tales como el análisis de datos mediante paneles de control y gráficos intuitivos, así como la capacidad de adaptarse a una variedad de contextos para los cuestionarios. Sin embargo, ninguna de estas herramientas realiza todas estas acciones de manera integral. Además, suelen omitir aspectos importantes, como la variabilidad de tipos de preguntas o la capacidad de crear cuestionarios dentro de la misma aplicación. Es importante destacar que todas estas aplicaciones están diseñadas para el ámbito de la educación superior. Es en este contexto que surge Kiddy Quiz, con el propósito de abordar todas estas áreas de manera integral, centrándose especialmente en los niños que cursan la educación básica elemental.
